Large language models can be used to extract candidate clinical
events from free-text electronic health records, but their outputs
require validation. Such a manual review produces labelled examples
that can subsequently be used to improve model performance. We
compare few-shot prompting, fine-tuning, and prompt tuning for
context-aware phrase classification, focusing on the identification
of drug discontinuation events in Estonian clinical notes. To gain
insights, we first examine how performance scales with training-set
size on a well-established public benchmark. We then validate the
methods on manually annotated candidate events from the Est-Health 30
dataset. We show that, in a low-data setting, prompt tuning matches
or outperforms full fine-tuning for instruction-tuned models while
training up to 10,000 times fewer parameters. These findings suggest
that soft prompts can effectively use limited clinical annotations
while requiring far fewer trainable parameters, although their
effectiveness depends on the underlying model.
